% Standalone section for inclusion in the main research report.
% Requires amsmath, amssymb, longtable, and hyperref in the enclosing document.
\section{Calibration of the compact K theory and the transferred machines}
\label{sec:k230-calibration}

\subsection{Pinned source and the exact consistency statement}

The reference in this section is Elliot Glazer's \texttt{k230-rocq}, commit
\path{3142e4dd980f0be12f4abb636c4ae2528afc17eb}.
Its literal table has 230 working states, start state 0, terminal label 230,
and SHA--256
\path{e7ccbcec074534cfd955f0ab94d16f13594e81974074309dc2ce663b9dfd7ba8}.
The reference files imported from Jeff's development are pinned to
\path{c132d7eae8158cd3a9850ea3eb7251158c3dd821}.
The retained source-manifest replay checks the bytes of all 216 proof
sources and all 460 transition rows. The same pinned reference was then
freshly compiled with the specified toolchain, its two final assumption
reports were closed, and an independent \texttt{coqchk} recheck succeeded.
The retained receipts distinguish source checks, compilation, and kernel
checking explicitly.

The symbol \(\mathrm{StandardZFC}\) is not an unspecified theory parameter.
In \texttt{KQuotZFC.v} it is defined by
\[
 \mathrm{StandardZFC}(A)\quad\Longleftrightarrow\quad
 (\mathrm{ZFc}(A)\lor A=\mathrm{Foundation})
       \lor A=\mathrm{Choice}.
\]
Here \texttt{Std\_iface.v} defines \(\mathrm{ZFc}\) to contain equality,
extensionality, empty set, pairing, union, power set, infinity, and all closed
instances of separation and replacement. The formulas use the ordinary
function-free membership-and-equality signature. The relation
\(T\vdash_{\mathrm C}\bot\) means a finite classical first-order derivation
of falsity from a finite list of axioms belonging to \(T\).
Consequently the consistency predicate at the endpoint is precisely
\(\neg(\mathrm{StandardZFC}\vdash_{\mathrm C}\bot)\).

\subsection{The single compact set-theoretic schema}

Let \(\varphi(y,p,\vec q)\) be an arbitrary formula, with the displayed
output and input variables and any finite tuple of additional parameters.
Use the abbreviations
\begin{align*}
 p\subseteq u&\;:\Longleftrightarrow\;\forall z\,(z\in p\to z\in u),\\
 W_\varphi(p,C)&\;:\Longleftrightarrow\;
   (\exists y\,\varphi(y,p,\vec q))\to
   \exists y\,(y\in C\land\varphi(y,p,\vec q)),\\
 H_\varphi(C)&\;:\Longleftrightarrow\;
   \forall u\bigl(u\in C\to
     ((\exists p\,\varphi(u,p,\vec q))\land
       \forall p\,(p\subseteq u\to W_\varphi(p,C)))\bigr).
\end{align*}
The ordinary theory \(K\), called \texttt{KPrimeOrd} in the development,
contains the equality axioms and the universal closures of
\begin{equation}
 \forall a\,\forall x\,\exists C\,
       ((\varphi(x,a,\vec q)\to x\in C)\land H_\varphi(C)).
 \label{eq:compact-k-schema}
\end{equation}
In particular, \(H_\varphi(C)\) remains required when
\(\varphi(x,a,\vec q)\) is false. Moving the entire closure condition under
that implication would change the theory and is not a permitted compression.

The emitted matrix uses a different local conjunction order and places a
support witness across a conjunct in which that witness variable does not
occur. These are proved correspondences in \texttt{KPrimeCore.v} and the
decoder bridge, rather than identifications made by reading a suggestive
formula. For numerical calibration the development also uses
\(K^{\mathrm{IR}}=K+\forall x\,\neg(x\in x)\), named
\texttt{KPrimeOrdIR}. The additional sentence is useful because the
distinguished numerical falsity code decodes to the membership atom
\(v_0\in v_0\). Both theories have the same consistency strength as
\(\mathrm{StandardZFC}\).

The two consistency directions are syntactic. For the lower reduction, raw
objects of \(K\) first provide a tagged domain and codes of relations.
A bisimulation quotient interprets extensional ZF without Foundation.
The power-set argument bounds canonical representatives and subsequently
includes all equivalent representatives; it does not assume that all raw
codes already form an extensional universe. The pinned reference's
\texttt{ac\_reg\_relcon} then supplies the reductions for Foundation and
Choice. For the upper reduction, constructible-hierarchy reflection provides
the required subset and witness closures. The proof collects the finitely
many ZFC axiom instances required for each finite \(K^{\mathrm{IR}}\)
derivation and translates that derivation into classical ZFC deduction.
Neither direction assumes that a model of ZFC exists.

\subsection{Numerical and machine interfaces}

The numerical constructor conventions are
\begin{align*}
 \mathrm{imp}(a,b)&=4P(a,b), &
 \mathrm{mem}(i,j)&=4P(i,j)+1,\\
 \mathrm{all}(i,a)&=4P(i,a)+2, &
 \mathrm{eq}(i,j)&=4P(i,j)+3.
\end{align*}
where \(P\) is the fixed Cantor pairing function. Falsity is code 1;
negation and conjunction are implication abbreviations relative to that
code. The axiom emitter retains all eight logical schemas, including both
membership congruence schemas, equality transitivity, and the quantifier
schema \((\exists a\,\forall a\,B)\to B\).
Its set-theory packet is
\(K(a,b,c)\land(V_{\in}(a,b,c)\land V_{=}(a,b,c))\), enabled only if
\(c\ne a,b\). The same number \(c\) serves as matrix code and private
variable label.

That last reuse requires proof at all code values. The lemmas
\texttt{D\_private\_fresh}, \texttt{D\_zero\_closed}, and
\texttt{D\_one} handle it, including codes 0 and 1. The separate proof
\texttt{named\_same\_roles} handles coincident public labels \(a=b\).
For completeness, \texttt{ordinary\_named\_core\_Q} uses a padded
ordinary formula code larger than the fixed public labels 1 and 3, so its
freshness gate is enabled. There is no assumption that every arbitrary
matrix code is an encoding of that specially padded form.

The actual blank bootstrap reaches the interpreter after 408419 steps,
with a gap of 181 zero cells, nine physical register values
\((0,0,0,3,0,0,3,0,0)\), and 167 spare registers containing 1.
The reference's source seed is therefore
\(\texttt{mkState}\;0\;3\;0\;0\;3\;0\), not an invented empty cache.
The complete chain is
\begin{align*}
 M_{230}\mathord\downarrow
 &\;\Longleftrightarrow\;\texttt{rm\_halts}\;p_{230}\;([],R_0)\\
 &\;\Longleftrightarrow\;\texttt{halts\_from}\;\texttt{seed230}\\
 &\;\Longleftrightarrow\;\texttt{thm230}\;1\\
 &\;\Longleftrightarrow\;K^{\mathrm{IR}}\vdash_{\mathrm C}\bot\\
 &\;\Longleftrightarrow\;\mathrm{StandardZFC}\vdash_{\mathrm C}\bot.
\end{align*}
The named interfaces are, in the same order:
\begin{quote}\small
\path{Sim230Main.tm230_rm}\\
\path{Main230.prog230}\\
\path{Seed230.seed230_iff_thm}\\
\path{Bridge230.thm230_iff_KPrimeOrdIR}\\
\path{KUpperConsistency.KPrimeOrdIR_standardZFC_equiconsistent}.
\end{quote}
\texttt{Correct230.v} composes them directly.

The simulation's symbolic register operations are not unexplained oracles.
The temporary section parameters \texttt{box\_mark}, \texttt{box\_inc},
\texttt{box\_dec0}, and \texttt{box\_decp} are discharged by the literal
tape proofs in \texttt{Tape230Blocks.v} and \texttt{Tape230Carry.v}.
\texttt{Interpreter230Check.v} checks every source leaf, and
\texttt{ProgramBisim230.v} checks both directions of the finite certificate
relating the literal interpreter tree to the structured source.

\subsection{Which earlier reductions transfer}

A native transformation that preserves the blank halting behavior of
\(M_{230}\) inherits the whole chain above. A source transformation that
preserves \texttt{rm\_halts}\,\(p_{230}\) from the actual seed inherits
the numerical and set-theoretic parts. This makes runtime and source
reductions composable with the stronger compact axiom package.

The following distinction is essential. Arbitrary reassociation of a
formula is not the same operation as reassociation of an interpreter
source tree: the former requires propositional derivability, while the
latter requires a control-flow or register-machine equivalence proof.
Closed source blocks can be rebracketed when their exit contracts are
respected and the resulting address-dependent runtime is reverified.
The old set-theory packet containing separate fixed ZF and collection
components cannot be substituted for the new \(K\) packet merely because
both have been described as ZFC-equiconsistent.

Logical-row permutations preserve the set of emitted schemas; their proof
compiler must select the new stack positions. Independent bijective
permutations of schematic parameters likewise preserve the set of axiom
instances. Finally, the antecedent exchange
\[
 (X\to(Y\to Z))\quad\longleftrightarrow\quad(Y\to(X\to Z))
\]
is available from the retained propositional basis:
\texttt{CD230.cd\_103} proves one implication and swapping \(X,Y\)
proves the other. These alternatives change neither equality strength nor
the quantifier rule. Removing either congruence axiom, or replacing the
quantifier schema by its superficially dual form, is not justified by
these transformations.

The closed increment-and-drain identity used here is proved in
Section~\ref{sec:k227-source-roles}. Its contextual replacement preserves
source halting for arbitrary natural register values. The changed addresses
still require the separate exact native decoder proof.

\subsection{External Rocq proof and the ZFC-internal assertion}

There are two different kinds of formalization here. The supplied Rocq
development is an external proof about finite syntax, finite derivations,
and executions. Its two final declarations state halting equivalence and
nonhalting equivalence. A successful \texttt{Print Assumptions} report
that they are closed under the global context rules out undischarged
theory or simulation assumptions in those declarations. A successful
\texttt{coqchk} recheck validates the compiled proof terms independently
of their producing tactics. The assumption claim concerns the dependency
closure of the named endpoints. Imported utility libraries also declare
unused axioms, which can appear in a module-wide context summary without
occurring in either endpoint proof. The pinned release uses Rocq 9.1.0 and
\texttt{rocq-stdlib} 9.0.0. Its documentation explicitly says that the
separate project concerning proofs internal to PA is not included.

The intended object-theoretic conclusion, for a transferred machine \(M\),
is
\begin{equation}
 \mathrm{ZFC}\vdash
 \bigl[\neg\exists t\,\mathrm{HaltTrace}_M(t)
       \;\longleftrightarrow\;
       \neg\exists p\,\mathrm{Proof}_{\mathrm{ZFC}}(p,\bot)\bigr].
 \label{eq:k-internal-target}
\end{equation}
The mathematical route to this internal assertion is finitary: encode
the fixed transition table, the register interpreter, formulas, proof
lists, and their translations by natural numbers. ZFC proves the elementary
recursion and induction needed to verify the interpreter simulations,
the proof enumerator's soundness and completeness, and the two syntactic
derivation translations. Applying those translations to arbitrary finite
proofs gives the two implications between the existential predicates in
\eqref{eq:k-internal-target}; classical negation then gives the displayed
equivalence. No appeal to ZFC's own consistency, its soundness, or existence
of a model of ZFC occurs in that argument.

The external kernel certificate and this internalization argument must not
be conflated. The release does not contain a serialized derivation of
\eqref{eq:k-internal-target} in an encoded ZFC proof calculus, and a fresh
kernel recheck alone would not create one. The defensible claims are the
exact external endpoint that was actually checked, together with the
finitary mathematical argument that its constructions can be formalized
in ZFC. Any later fully internal checker can target the same fixed
arithmetized statement without changing the machine.
